\chapter{Themenregister über alle Bände}

Wo finde ich was? Die folgende Übersicht ordnet die wichtigsten Begriffe und Sätze den Bänden zu, in denen sie behandelt werden. Zuerst genannt ist jeweils der Band mit der ausführlichsten Darstellung.

\section*{Grundbegriffe}
\addcontentsline{toc}{section}{Grundbegriffe}
\nopagebreak
{\small\renewcommand{\arraystretch}{1.32}
\begin{longtable}{@{}P{5.6cm}P{2.8cm}@{}}
\toprule
\textbf{Begriff} & \textbf{Bände}\\
\midrule
\endhead
Zahlbereiche N, Z, Q, R, C & 0, 1, 3\\
Irrationalität von 2 & 0, 1, 7\\
Bruchrechnung & 0\\
Potenzen, Wurzeln, Logarithmen & 0\\
Binomische Formeln & 0, 6\\
Quadratische Gleichungen & 0, 7\\
Ungleichungen und Beträge & 0, 3\\
Funktionsbegriff & 0, 3, 7, 10\\
Umkehrfunktion, Verkettung & 0, 3\\
Vollständige Induktion & 0, 1, B\\
Mengen, Relationen, Abbildungen & 10, B\\
\bottomrule
\end{longtable}}

\section*{Analysis}
\addcontentsline{toc}{section}{Analysis}
\nopagebreak
{\small\renewcommand{\arraystretch}{1.32}
\begin{longtable}{@{}P{5.6cm}P{2.8cm}@{}}
\toprule
\textbf{Begriff} & \textbf{Bände}\\
\midrule
\endhead
Folgen und Grenzwerte & 0, 3\\
$\varepsilon$-$\delta$-Definition & 0, 3, 7, 10\\
Reihen, geometrische Reihe & 0, 3\\
Stetigkeit, Zwischenwertsatz & 3, 8\\
Gleichmäßige Stetigkeit & 3, 10\\
Ableitung, Ableitungsregeln & 0, 3\\
Mittelwertsätze, Taylor & 3\\
Kurvendiskussion, Extremwerte & 0, 3\\
Riemann-Integral & 0, 3\\
Hauptsatz der Analysis & 0, 3, 7\\
Integrationstechniken & 3, 5\\
Gleichmäßige Konvergenz & 3\\
Differentialrechnung im $\R^{n}$ & 3, 9\\
Kurven- und Flächenintegrale & 9\\
Sätze von Green, Gauß, Stokes & 9\\
Differentialformen & 9\\
Nichtstandard-Analysis & 3, 10\\
Konvexität, Jensensche Ungleichung & 19, 22, 24, 25, 20, 3\\
Lagrange-Multiplikatoren & 19, 3\\
Logistische Gleichung & 18, 14, 21\\
\bottomrule
\end{longtable}}

\section*{Algebra und Zahlentheorie}
\addcontentsline{toc}{section}{Algebra und Zahlentheorie}
\nopagebreak
{\small\renewcommand{\arraystretch}{1.32}
\begin{longtable}{@{}P{5.6cm}P{2.8cm}@{}}
\toprule
\textbf{Begriff} & \textbf{Bände}\\
\midrule
\endhead
Teilbarkeit, ggT, euklidischer Algorithmus & 1\\
Primzahlen, Fundamentalsatz & 1, 7\\
Unendlichkeit der Primzahlen & 0, 1, B\\
Modulares Rechnen, Kongruenzen & 1\\
Chinesischer Restsatz & 1\\
Sätze von Fermat, Euler, Wilson & 1\\
Quadratische Reziprozität & 1, 7\\
Primzahlverteilung, Riemannsche Vermutung & 1\\
Diophantische Gleichungen & 1, 10\\
Kettenbrüche, Pellsche Gleichung & 1\\
Elliptische Kurven & 1\\
Kryptographie, RSA & 1\\
Lineare Gleichungssysteme & 0, 2, 5\\
Matrizen, Determinante & 2\\
Vektorräume, Basis, Dimension & 2\\
Lineare Abbildungen & 2\\
Eigenwerte, Jordanform & 2, 5\\
Skalarprodukt, Spektralsatz & 2, 9\\
Singulärwertzerlegung & 2, 5\\
Gruppen und Symmetrie & 6, 8\\
Tensoren, Tensoralgebra & 9\\
\bottomrule
\end{longtable}}

\section*{Geometrie und Topologie}
\addcontentsline{toc}{section}{Geometrie und Topologie}
\nopagebreak
{\small\renewcommand{\arraystretch}{1.32}
\begin{longtable}{@{}P{5.6cm}P{2.8cm}@{}}
\toprule
\textbf{Begriff} & \textbf{Bände}\\
\midrule
\endhead
Satz des Pythagoras & 0, 6, 7\\
Strahlensätze, Ähnlichkeit & 0, 6\\
Flächen und Volumina & 0, 6\\
Trigonometrie, Einheitskreis & 0, 6\\
Sinus- und Kosinussatz & 0, 6\\
Euklids Axiome & 6, 7, B\\
Kegelschnitte & 6\\
Isometrien, Symmetriegruppen & 6, 8\\
Projektive Geometrie & 6\\
Parallelenproblem & 6, 7, B\\
Hyperbolische, sphärische Geometrie & 6, 7\\
Krümmung, Gauß-Bonnet & 6, 9\\
Metrische und topologische Räume & 8\\
Kompaktheit, Zusammenhang & 8, 10\\
Fundamentalgruppe, Homotopie & 8\\
Klassifikation der Flächen & 8\\
Homologie & 8\\
Knotentheorie & 8\\
Mannigfaltigkeiten & 8, 9\\
Poincaré-Vermutung & 8\\
\bottomrule
\end{longtable}}

\section*{Höhere Algebra und Funktionentheorie}
\addcontentsline{toc}{section}{Höhere Algebra und Funktionentheorie}
\nopagebreak
{\small\renewcommand{\arraystretch}{1.32}
\begin{longtable}{@{}P{5.6cm}P{2.8cm}@{}}
\toprule
\textbf{Begriff} & \textbf{Bände}\\
\midrule
\endhead
Gruppe, Untergruppe, Ordnung & 11, 6, 15\\
Satz von Lagrange & 11, 1\\
Normalteiler, Faktorgruppe & 11\\
Homomorphiesätze & 11, 2\\
Gruppenoperation, Bahnensatz & 11, 15\\
Sylowsätze & 11\\
Ring, Ideal, Hauptidealring & 11, 1\\
Polynomring, Irreduzibilität & 11, 0\\
Körpererweiterung, Grad & 11\\
Konstruierbarkeit mit Zirkel und Lineal & 11, 6, 7\\
Endliche Körper & 11, 1, 15\\
Galoiskorrespondenz & 11, 7\\
Unlösbarkeit der Quintik & 11, 7\\
Komplexe Zahlen, Polarform & 12, 0, 2\\
Holomorphie, Cauchy-Riemann & 12, 3\\
Cauchyscher Integralsatz & 12, 9\\
Potenz- und Laurentreihen & 12, 3\\
Residuensatz & 12\\
Konforme Abbildungen & 12, 6, 8\\
Riemannscher Abbildungssatz & 12, 13\\
Gammafunktion & 12, 15\\
Zetafunktion, Riemannsche Vermutung & 12, 1\\
Defekte des Riemann-Integrals & 13, 3\\
$\sigma$-Algebra, Maß & 13, 4\\
Menge von Vitali, Banach-Tarski & 13, 10\\
Lebesgue-Integral & 13, 3, 4\\
Konvergenzsätze (monoton, dominiert) & 13\\
Fubini, Produktmaße & 13, 3\\
$L^{p}$-Räume, Hölder, Minkowski & 13\\
Banachraum, Hilbertraum & 13, 2\\
Hahn-Banach & 13\\
Banach-Steinhaus, offene Abbildung & 13\\
Orthonormalbasis, Fourierreihe & 13, 3, 14\\
Kompakte Operatoren, Spektralsatz & 13, 2, 14\\
Sobolevraum, schwache Ableitung & 13, 14\\
Banachscher Fixpunktsatz & 13, 14, 5\\
Richtungsfeld, Anfangswertproblem & 14\\
Trennung der Variablen & 14, 0\\
Picard-Lindelöf, Peano & 14\\
Lemma von Gronwall & 14\\
Matrixexponential & 14, 2\\
Resonanz, harmonischer Oszillator & 14, 9\\
Besselfunktionen, Legendrepolynome & 14, 12\\
Sturm-Liouville, Eigenwertproblem & 14, 13\\
Phasenporträt, Stabilität, Lyapunov & 14\\
Bifurkation, Grenzzyklus, Chaos & 14, 8\\
Wellengleichung, d’Alembert & 14, 9\\
Wärmeleitungsgleichung & 14, 4, 13\\
Laplacegleichung, harmonische Funktionen & 14, 12, 9\\
Lax-Milgram, finite Elemente & 14, 5\\
Distributionen, Deltadistribution & 14, 13\\
\bottomrule
\end{longtable}}

\section*{Stochastik und Numerik}
\addcontentsline{toc}{section}{Stochastik und Numerik}
\nopagebreak
{\small\renewcommand{\arraystretch}{1.32}
\begin{longtable}{@{}P{5.6cm}P{2.8cm}@{}}
\toprule
\textbf{Begriff} & \textbf{Bände}\\
\midrule
\endhead
Kombinatorik & 4\\
Bedingte Wahrscheinlichkeit, Bayes & 4\\
Zufallsvariablen, Verteilungen & 4\\
Erwartungswert, Varianz & 4\\
Gesetze der großen Zahlen & 4\\
Zentraler Grenzwertsatz & 4\\
Markow-Ketten, Martingale & 4\\
Schätzen und Testen & 17, 22, 23, 25, 26, 4\\
Gleitkommazahlen, Rundungsfehler & 5\\
Kondition und Stabilität & 5\\
LRund QR-Zerlegung & 5\\
Newton-Verfahren & 5\\
Interpolation, Splines & 5\\
Numerische Quadratur & 5\\
Differentialgleichungen, numerisch & 5\\
Monte-Carlo-Verfahren & 4, 5\\
Konfidenzintervalle & 17, 4\\
Fehler erster und zweiter Art & 17\\
Fallzahlplanung & 17, 20\\
Lineare Regression & 17, 22, 23, 2\\
Logistische Regression & 22, 17, 23\\
Überlebenszeitanalyse & 17\\
Bayessche Verfahren & 17, 4\\
Multiples Testen & 17\\
Metaanalyse & 17\\
Replikationskrise & 17, 18\\
Fisher-Information, Cramér-Rao & 20, 17\\
Relative Entropie, Kullback-Leibler & 20, 22, 17, 21\\
Ordnungsstatistiken & 21, 4\\
\bottomrule
\end{longtable}}

\section*{Kombinatorik und Graphentheorie}
\addcontentsline{toc}{section}{Kombinatorik und Graphentheorie}
\nopagebreak
{\small\renewcommand{\arraystretch}{1.32}
\begin{longtable}{@{}P{5.6cm}P{2.8cm}@{}}
\toprule
\textbf{Begriff} & \textbf{Bände}\\
\midrule
\endhead
Summen- und Produktregel & 15, 4\\
Bijektionsprinzip & 15, 10\\
Doppeltes Abzählen & 15\\
Schubfachprinzip & 15, 1\\
Fakultät, Stirlingformel & 15, 4, 12\\
Binomialkoeffizient, Pascalsches Dreieck & 15, 0, 4\\
Sterne und Balken, Multimengen & 15\\
Vandermonde, Hockeyschläger & 15\\
Satz von Lucas & 15, 1\\
Siebformel, Derangements & 15, 4\\
Erzeugende Funktionen & 15, 4, 12\\
Catalanzahlen & 15\\
Stirlingzahlen, Bellzahlen & 15\\
Partitionen, Ferrersdiagramm & 15, 1\\
Zwölffachweg & 15\\
Burnside, Pólya, Zyklenzeiger & 15, 11\\
Graph, Grad, Handschlaglemma & 15, 8\\
Bäume, Cayley-Formel, Prüfercode & 15, 5\\
Eulerkreis, Hamiltonkreis & 16, 15, 7\\
Planarität, Polyederformel, Kuratowski & 15, 6, 8\\
Chromatische Zahl, Vierfarbensatz & 15, 8, B\\
Matching, Satz von Hall & 15\\
Max-Flow-Min-Cut, Menger & 15, 5\\
Ramseyzahlen, van der Waerden, Schur & 15\\
Probabilistische Methode & 15, 4\\
Turán, Sperner, LYM & 15\\
Lovász Local Lemma & 15\\
Zufallsgraphen, Schwellenwerte & 16, 15, 4\\
Handschlaglemma, Gradfolgen & 16, 15\\
Spannbäume, Prüfercode & 16, 15\\
Planarität, Vierfarbensatz & 16, 15, B\\
Färbungen, chromatische Zahl & 16, 15\\
Matchings, Satz von Hall & 16, 15, 13\\
Flüsse, Max-Flow-Min-Cut & 16, 19, 15\\
Zentralitätsmaße, kleine Welten & 16\\
Ausbreitung auf Netzwerken & 16, 18\\
\bottomrule
\end{longtable}}

\section*{Optimierung und Entscheidung}
\addcontentsline{toc}{section}{Optimierung und Entscheidung}
\nopagebreak
{\small\renewcommand{\arraystretch}{1.32}
\begin{longtable}{@{}P{5.6cm}P{2.8cm}@{}}
\toprule
\textbf{Begriff} & \textbf{Bände}\\
\midrule
\endhead
Lineare Programme, Simplexverfahren & 19, 5\\
Dualität, Schattenpreise & 19, 21\\
Komplementärer Schlupf & 19, 20\\
Totale Unimodularität & 19, 16\\
Karush-Kuhn-Tucker & 19, 3\\
Gradienten- und Newtonverfahren & 19, 22, 5\\
Innere-Punkte-Verfahren & 19\\
Ganzzahlige Optimierung & 19, 16\\
Verzweigen und Beschränken & 19\\
P gegen NP, NP-Vollständigkeit & 19, 10, B\\
Robuste Optimierung & 19, 21\\
Pareto-Front, mehrere Ziele & 19, 21\\
\bottomrule
\end{longtable}}

\section*{Information und Kryptographie}
\addcontentsline{toc}{section}{Information und Kryptographie}
\nopagebreak
{\small\renewcommand{\arraystretch}{1.32}
\begin{longtable}{@{}P{5.6cm}P{2.8cm}@{}}
\toprule
\textbf{Begriff} & \textbf{Bände}\\
\midrule
\endhead
Informationsgehalt, Bit & 20\\
Entropie, bedingte Entropie & 20, 22, 4\\
Quellencodierung, Huffman & 20\\
Kraftsche Ungleichung & 20\\
Transinformation, Kanalkapazität & 20\\
Kanalcodierungstheorem & 20\\
Typische Mengen & 20, 4\\
Fehlerkorrigierende Codes & 20, 16\\
Kolmogorow-Komplexität & 20, 10, B\\
Kürzeste Beschreibung, MDL & 20, 22, 17\\
Perfekte Sicherheit, Einmalschlüssel & 20\\
Blockchiffren, Betriebsarten & 20\\
Hashfunktionen, Geburtstagsschranke & 20, 15\\
Diffie-Hellman, diskreter Logarithmus & 20, 1, 11\\
RSA und Faktorisierung & 20, 1\\
Elliptische Kurven & 20, 1\\
Post-Quanten-Verfahren & 20\\
\bottomrule
\end{longtable}}

\section*{Spiele und Strategien}
\addcontentsline{toc}{section}{Spiele und Strategien}
\nopagebreak
{\small\renewcommand{\arraystretch}{1.32}
\begin{longtable}{@{}P{5.6cm}P{2.8cm}@{}}
\toprule
\textbf{Begriff} & \textbf{Bände}\\
\midrule
\endhead
Normalform, Dominanz & 21\\
Sattelpunkt, Maximin und Minimax & 21\\
Minimax-Satz & 21, 19\\
Gemischte Strategien & 21, 4\\
Nash-Gleichgewicht & 21\\
Nutzentheorie nach von Neumann & 21, 4\\
Gefangenendilemma & 21, 11\\
Feigheitsspiel, Hirschjagd & 21\\
Fixpunktsatz von Brouwer & 21, 8, 14\\
Extensive Form, Teilspielperfektheit & 21\\
Wiederholte Spiele, Volkssatz & 21\\
Evolutionär stabile Strategien & 21, 18\\
Replikatordynamik & 21, 18, 14\\
Kollektivgüter, Allmendetragik & 21\\
Satz von Arrow & 21\\
Auktionen, Zweitpreisregel & 21\\
Stabile Zuordnung, Gale-Shapley & 21, 16\\
Kern und Shapley-Wert & 21\\
\bottomrule
\end{longtable}}

\section*{Modelle in Biologie und Medizin}
\addcontentsline{toc}{section}{Modelle in Biologie und Medizin}
\nopagebreak
{\small\renewcommand{\arraystretch}{1.32}
\begin{longtable}{@{}P{5.6cm}P{2.8cm}@{}}
\toprule
\textbf{Begriff} & \textbf{Bände}\\
\midrule
\endhead
Exponentielles und logistisches Wachstum & 18, 14\\
Chaos in der Rekursion & 18, 14\\
Lotka-Volterra, Wettbewerb & 18, 14, 21\\
Stabilität, Phasenporträt & 18, 14\\
SIR-Modell, Basisreproduktionszahl & 18\\
Impfschwelle & 18, 21\\
Pharmakokinetik, Kompartimente & 18, 14\\
Michaelis-Menten-Kinetik & 18\\
FitzHugh-Nagumo, erregbare Systeme & 18, 14\\
Gompertz-Wachstum, Tumormodelle & 18\\
Resistenz, adaptive Therapie & 18, 21\\
Goldie-Coldman, Luria-Delbrück & 18, 4\\
\bottomrule
\end{longtable}}

\section*{Lernen aus Daten}
\addcontentsline{toc}{section}{Lernen aus Daten}
\nopagebreak
{\small\renewcommand{\arraystretch}{1.32}
\begin{longtable}{@{}P{5.6cm}P{2.8cm}@{}}
\toprule
\textbf{Begriff} & \textbf{Bände}\\
\midrule
\endhead
Bayes-Regel, Bayes-Risiko & 22, 4\\
Verzerrung und Varianz & 22, 17\\
Überanpassung, Regularisierung & 22, 17, 23, 5\\
Kreuzvalidierung, Datenleck & 22, 17\\
Trennschärfe, Kalibrierung & 22, 17\\
Kammregression, Lasso & 22, 19\\
Perzeptron, Stützvektormaschine & 22, 19\\
Kerntrick & 22, 6\\
PAC-Lernen, VC-Dimension & 22, 10\\
Verallgemeinerungsschranken & 22\\
Entscheidungsbäume, Zufallswald & 22, 23, 16\\
Boosting & 22\\
Nächste Nachbarn & 22\\
Fluch der Dimension & 22, 19\\
Neuronale Netze, Rückpropagation & 22\\
Universelle Näherung & 22, 13, 3\\
Clusterbildung, EM-Verfahren & 22, 17\\
Hauptkomponenten & 22, 2, 17\\
Doppelter Abstieg & 22\\
Implizite Regularisierung & 22\\
\bottomrule
\end{longtable}}

\section*{Ursache und Wirkung}
\addcontentsline{toc}{section}{Ursache und Wirkung}
\nopagebreak
{\small\renewcommand{\arraystretch}{1.32}
\begin{longtable}{@{}P{5.6cm}P{2.8cm}@{}}
\toprule
\textbf{Begriff} & \textbf{Bände}\\
\midrule
\endhead
Simpsons Paradoxon & 23, 17\\
Potentielle Ergebnisse & 23\\
Randomisierung, Versuchsplanung & 23, 17\\
Kausale Graphen, d-Trennung & 23, 16\\
Störgröße, Mediator, Kollider & 23\\
Hintertür- und Vordertürkriterium & 23\\
Behandlungsneigung & 23\\
Gewichtung nach Kehrwahrscheinlichkeit & 23, 17\\
Doppelt robuste Schätzung & 23, 22\\
Instrumentvariablen & 23\\
Regressionssprung, doppelte Differenz & 23\\
Zielstudienemulation & 23\\
g-Formel, zeitabhängige Störgrößen & 23\\
Mediation, direkte und indirekte Effekte & 23\\
Heterogene Behandlungseffekte & 23, 22\\
Empfindlichkeitsanalyse & 23, 17\\
Unsterbliche Zeit & 23\\
\bottomrule
\end{longtable}}

\section*{Entscheiden und Bewerten}
\addcontentsline{toc}{section}{Entscheiden und Bewerten}
\nopagebreak
{\small\renewcommand{\arraystretch}{1.32}
\begin{longtable}{@{}P{5.6cm}P{2.8cm}@{}}
\toprule
\textbf{Begriff} & \textbf{Bände}\\
\midrule
\endhead
Erwartungsnutzen, Nutzenaxiome & 24, 21\\
Risikoscheu, Sicherheitsäquivalent & 24\\
Behandlungs- und Testschwelle & 24\\
Wert der Information & 24, 20\\
Entscheidungsbaum & 24, 22\\
Empfindlichkeitsanalyse & 24, 23\\
Markow-Modell, QALY & 24, 18\\
Diskontierung & 24\\
Dynamische Programmierung, Bellman & 24, 19\\
Markow-Entscheidungsprozess & 24\\
Bandit, Erkunden gegen Ausnutzen & 24, 22\\
Kosten-Wirksamkeit, ICER & 24\\
Ambiguität, Minimax & 24, 21\\
Neue Erwartungstheorie, Rahmung & 24\\
Satz von Arrow & 24, 21\\
Priorisierung, Verteilungsgerechtigkeit & 24\\
\bottomrule
\end{longtable}}

\section*{Prozesse und Zufall über die Zeit}
\addcontentsline{toc}{section}{Prozesse und Zufall über die Zeit}
\nopagebreak
{\small\renewcommand{\arraystretch}{1.32}
\begin{longtable}{@{}P{5.6cm}P{2.8cm}@{}}
\toprule
\textbf{Begriff} & \textbf{Bände}\\
\midrule
\endhead
Irrfahrt, Arkussinus-Gesetz & 25, 16\\
Markow-Kette, Übergangsmatrix & 25, 24, 4\\
Stationäre Verteilung, Mischzeit & 25, 16\\
Absorption, Fundamentalmatrix & 25, 24\\
Poisson-Prozess, Wartezeiten & 25, 4\\
Warteschlangen, Little-Formel & 25, 19\\
Verzweigungsprozess, Aussterben & 25, 18\\
Bedingte Erwartung, Filtration & 25, 17, 23, 26\\
Martingal, optionales Stoppen & 25\\
Sequenzielles Testen & 25, 17\\
Brownsche Bewegung & 25\\
Itō-Integral, Itō-Formel & 25\\
Stochastische Differentialgleichung & 25, 14\\
Optionsbewertung, Black-Scholes & 25\\
Genetische Drift, Fixierung & 25, 18\\
Kalman-Filter & 25, 22\\
\bottomrule
\end{longtable}}

\section*{Ökonometrie und Beobachtungsdaten}
\addcontentsline{toc}{section}{Ökonometrie und Beobachtungsdaten}
\nopagebreak
{\small\renewcommand{\arraystretch}{1.32}
\begin{longtable}{@{}P{5.6cm}P{2.8cm}@{}}
\toprule
\textbf{Begriff} & \textbf{Bände}\\
\midrule
\endhead
Endogenität, Simultanität & 26, 23\\
Identifikation, Abzählbedingung & 26, 23\\
Instrumentvariablen, schwache Instrumente & 26, 23\\
Momentenmethode, GMM & 26, 17\\
Logit, Probit, diskrete Wahl & 26, 17, 24\\
Zensierung, Tobit, Zweiteilmodell & 26\\
Auswahlverzerrung, Heckman & 26, 23\\
Paneldaten, feste Effekte & 26\\
Dynamische Panels, Nickell & 26\\
Doppelte Differenz, parallele Trends & 26, 23\\
Stationarität, Einheitswurzel & 26, 25\\
Scheinregression & 26, 25\\
Kointegration, Fehlerkorrektur & 26\\
Vektorautoregression, Granger-Kausalität & 26\\
Lucas-Kritik, reduzierte Form & 26, 21\\
Vorhersage gegen Politik & 26, 22, 23\\
\bottomrule
\end{longtable}}

\section*{Simulation im Rechner}
\addcontentsline{toc}{section}{Simulation im Rechner}
\nopagebreak
{\small\renewcommand{\arraystretch}{1.32}
\begin{longtable}{@{}P{5.6cm}P{2.8cm}@{}}
\toprule
\textbf{Begriff} & \textbf{Bände}\\
\midrule
\endhead
Pseudozufallszahlen, lineare Kongruenz & 27, 4\\
Inversion, Verwerfung, Box--Muller & 27, 4\\
Monte-Carlo-Integration & 27, 5\\
Fluch der Dimension & 27, 22, 5\\
Standardfehler bei Abhängigkeit, Blockmittel & 27, 17\\
Antithetische Variable, Kontrollvariable, Schichtung & 27\\
Wichtigkeitsstichprobe & 27, 25\\
Metropolis, Gibbs-Sampler & 27, 25, 22\\
Konvergenzdiagnostik, $\hat R$, wirksames $n$ & 27\\
Hamilton-Monte-Carlo & 27\\
Euler--Maruyama, Milstein, Mehrebenen-Monte-Carlo & 27, 25\\
Ereignisgesteuerte Simulation & 27, 25\\
Seltene Ereignisse, Aufspaltung & 27\\
Teilchenfilter & 27, 25\\
Reproduzierbarkeit, Startwert & 27\\
Unsicherheit des Simulationsergebnisses & 27, 17\\
\bottomrule
\end{longtable}}
\section*{Dynamische Systeme und Chaos}
\addcontentsline{toc}{section}{Dynamische Systeme und Chaos}
\nopagebreak
{\small\renewcommand{\arraystretch}{1.32}
\begin{longtable}{@{}P{5.6cm}P{2.8cm}@{}}
\toprule
\textbf{Begriff} & \textbf{Bände}\\
\midrule
\endhead
Fluss, Phasenraum, Fixpunkt & 28, 14\\
Linearisierung, Stabilität, Eigenwerte & 28, 14, 2\\
Grenzzyklus, Poincaré--Bendixson & 28, 14\\
Sattel-Knoten-, Heugabel-, Hopf-Bifurkation & 28, 14\\
Erzwungene Schwingung, Synchronisation & 28\\
Periodenverdopplung, Feigenbaum & 28\\
Lorenz-System, seltsamer Attraktor & 28, 14\\
Lyapunov-Exponent, Vorhersagehorizont & 28\\
Einbettung, Satz von Takens & 28\\
Fraktale Dimension, Korrelationsdimension & 28, 8\\
Surrogatdaten, Test gegen Zufall & 28, 17\\
Reaktions-Diffusions-Muster, Spiralwelle & 28, 18\\
Gekoppelte Oszillatoren, Kuramoto & 28\\
Erregbare Zelle, Schrittmacher, Flimmern & 28, 18\\
\bottomrule
\end{longtable}}
\section*{Algebraische Geometrie}
\addcontentsline{toc}{section}{Algebraische Geometrie}
\nopagebreak
{\small\renewcommand{\arraystretch}{1.32}
\begin{longtable}{@{}P{5.6cm}P{2.8cm}@{}}
\toprule
\textbf{Begriff} & \textbf{Bände}\\
\midrule
\endhead
Affine Varietät, Zariski-Topologie & 29\\
Ideal, Verschwindungsideal, Radikal & 29, 11\\
Hilberts Basissatz & 29, 11\\
Nullstellensatz & 29\\
Monomordnung, Division mit Rest & 29\\
Gröbner-Basis, Buchberger-Algorithmus & 29\\
Elimination, Resultante, Implizitisierung & 29, 11\\
Hilbert-Funktion, Dimension, Grad & 29\\
Tangentialraum, Glattheit, Jacobi-Kriterium & 29, 9\\
Projektiver Raum, Satz von Bézout & 29, 6\\
Elliptische Kurve, Gruppengesetz & 29, 1, 20\\
Rationale Parametrisierung, Lüroth & 29\\
Geschlecht einer Kurve & 29, 12\\
Identifizierbarkeit, algebraische Statistik & 29, 18, 23\\
Inverse Kinematik, Kurven in der Kryptographie & 29, 20\\
\bottomrule
\end{longtable}}
\section*{Darstellungstheorie}
\addcontentsline{toc}{section}{Darstellungstheorie}
\nopagebreak
{\small\renewcommand{\arraystretch}{1.32}
\begin{longtable}{@{}P{5.6cm}P{2.8cm}@{}}
\toprule
\textbf{Begriff} & \textbf{Bände}\\
\midrule
\endhead
Punktgruppe, Symmetrieoperation, Schoenflies & 30, 11, 6\\
Darstellung, Unitarität, invarianter Unterraum & 30, 2\\
Satz von Maschke, Schurs Lemma & 30, 11\\
Charakter, Orthogonalitätsrelationen & 30\\
Charaktertafel, Mulliken-Symbole & 30\\
Normalschwingung, Auswahlregel, Alternativverbot & 30, 14\\
Projektionsformel, Symmetriekoordinaten & 30\\
Einschränkung, Induktion, Frobenius-Reziprozität & 30\\
Symmetrische Gruppe, Young-Diagramm, Hakenformel & 30, 15, 11\\
Tensorprodukt, Clebsch--Gordan, Frobenius--Schur-Indikator & 30\\
Haar-Maß, kompakte Gruppe, Peter--Weyl & 30, 13\\
$SO(3)$, Kugelflächenfunktionen, Orbitale & 30, 9\\
$SU(2)$, Spin, Kramers & 30\\
Fourier auf Gruppen, Schur--Weyl & 30, 12\\
Hückel-Orbitale, Aromatizität & 30\\
Kristallfeld, high-spin/low-spin & 30\\
Raumgruppen, Bloch, Quasikristalle & 30, 6\\
$SU(3)$, Quarks, Achtfacher Weg & 30\\
\bottomrule
\end{longtable}}

\section*{Logik, Grundlagen, Geschichte}
\addcontentsline{toc}{section}{Logik, Grundlagen, Geschichte}
\nopagebreak
{\small\renewcommand{\arraystretch}{1.32}
\begin{longtable}{@{}P{5.6cm}P{2.8cm}@{}}
\toprule
\textbf{Begriff} & \textbf{Bände}\\
\midrule
\endhead
Beweistechniken (direkt, indirekt, Induktion) & 0, B\\
Aussagen- und Prädikatenlogik & 10, B\\
Beweiskalküle, natürliches Schließen & 10, B\\
Gödelscher Vollständigkeitssatz & 10, B\\
Kompaktheitssatz & 10\\
Löwenheim-Skolem, Skolem-Paradoxon & 10\\
Russellsche Antinomie & 10, 7, B\\
ZFC-Axiome & 10, B\\
Ordinal- und Kardinalzahlen & 10\\
Satz von Cantor, Diagonalverfahren & 3, 10, B\\
Auswahlaxiom, Zornsches Lemma & 10, 2\\
Banach-Tarski-Paradoxon & 10\\
Berechenbarkeit, Halteproblem & 10, B\\
Gödels Unvollständigkeitssätze & 10, B, 7\\
Satz von Tarski & 10\\
Modelltheorie, Ultraprodukte & 10\\
Forcing, Kontinuumshypothese & 10, B\\
Große Kardinalzahlen & 10, B\\
Typentheorie, HoTT & 10, B\\
Beweisassistenten, Formalisierung & 10, B\\
Grundlagenkrise & 7, B\\
Hilberts Programm & 7, 10, B\\
Geschichte der Zahlzeichen & 7, 0\\
Griechische Mathematik & 7, 6, B\\
Islamische Mathematik & 7\\
Infinitesimalrechnung, Prioritätsstreit & 7, 0, 3\\
Rigorisierung des 19. Jahrhunderts & 7, 3\\
\bottomrule
\end{longtable}}

\section*{Die großen offenen Fragen}
\addcontentsline{toc}{section}{Die großen offenen Fragen}
\nopagebreak
{\small\renewcommand{\arraystretch}{1.32}
\begin{longtable}{@{}P{5.0cm}P{1.6cm}P{1.8cm}@{}}
\toprule
\textbf{Frage} & \textbf{seit} & \textbf{Bände}\\
\midrule
\endhead
Riemannsche Vermutung & 1859 & 1, 12\\
Goldbachsche Vermutung & 1742 & 1\\
Primzahlzwillinge & Antike & 1\\
P gegen NP & 1971 & 19, 10, B\\
Navier-Stokes-Regularität & 1934 & 14, 3, 9\\
Birch und Swinnerton-Dyer & 1965 & 1\\
Hodge-Vermutung & 1950 & 8\\
Yang-Mills-Massenlücke & 1954 & 9\\
Kontinuumshypothese (Status) & 1878 & 10, B\\
Existenz des ultimativen L & 2010 & 10\\
Hilberts 10. Problem über $\mathbb{Q}$ & 1900 & 10\\
Kugelpackungen in Dimension $n$ & 1611 & 6\\
Untere Schranken der Beweiskomplexität & 1985 & 10, B\\
Hilberts 16. Problem (Grenzzyklen) & 1900 & 14\\
Ramseyzahl $R(5,5)$ & 1955 & 15, 16\\
Asymptotik von $R(s,s)$ & 1935 & 15, 16\\
Union-Closed-Sets-Vermutung & 1979 & 15\\
Sonnenblumenvermutung & 1960 & 15\\
Happy-End-Problem $\mathrm{ES}(n)$ & 1935 & 15, 6\\
Hadwiger-Vermutung & 1943 & 16, 15\\
Rekonstruktionsvermutung & 1942 & 16\\
Kapazität des Interferenzkanals & 1970er & 20\\
Sicherheit der Gitterverfahren & 2000er & 20\\
$\mathsf{PPAD}$ gegen $\mathsf{P}$ & 2006 & 21\\
Auswahl unter mehreren Gleichgewichten & 1950er & 21\\
Verallgemeinerung tiefer Netze & 2017 & 22\\
Gutartige Überanpassung & 2019 & 22\\
Herkunft der Skalengesetze & 2020 & 22\\
Übertragbarkeit kausaler Effekte & 2010er & 23\\
Struktur aus Daten gewinnen & 1990er & 23\\
Nutzen ohne Beeinflussung erheben & 1980er & 24\\
Gerechtigkeit in der Priorisierung & 1990er & 24\\
Selbstvermeidende Irrfahrt & 1949 & 25\\
Mischzeit im Einzelfall & 1980er & 25\\
Externe Gültigkeit lokaler Effekte & 2010er & 26\\
Erwartungen messen statt unterstellen & 1976 & 26\\
Schnelle Fourier-Transformation auf allen Gruppen & 1989 & 30, 27\\
Kronecker-Koeffizienten & 1938 & 30\\
Lorenz-Attraktor: Dimension exakt & 1963 & 28\\
Hilberts 16. Problem, zweiter Teil & 1900 & 28, 14\\
Minimale Gröbner-Komplexität im Einzelfall & 1982 & 29\\
Rang elliptischer Kurven, beschränkt? & 1922 & 29, 1\\
\bottomrule
\end{longtable}}

\bigskip
\begin{flushright}
\itshape\color{grau}
Nichts ist selbstverständlich.\\
Nichts ist trivial.\\
Jeder Schritt wird erklärt.\\[10pt]
\normalfont\color{grau}\upshape
--- die Regel, nach der\\
\itshape jede Seite geschrieben ist.
\end{flushright}